\documentclass{article}
\usepackage[en]{homework}

\config{Analysis-0}{2026 Spring}{6}

\begin{document}


\begin{problem}[3-8]
If $\sum a_n$ converges and $\{b_n\}$ is monotone and bounded, prove that
\[
\sum a_n b_n
\]
also converges.
\end{problem}

\begin{proof}
  Denote $s_n = a_0 + a_1 + \cdots + a_n$. By the Abel summation formula, we have
  \[ S_n:=\sum_{k=1}^{n}a_kb_k=s_nb_n+\sum_{k=1}^{n-1}s_k(b_k-b_{k+1}). \] \par
  We will prove that \( \sum a_nb_n \) is a Cauchy series, hence \( \sum a_nb_n \) converges. Note that \( \left\{s_n\right\} \) and \( \left(b_n\right) \) all converge, thus \( \left\{s_nb_n\right\} \) also converges, and there exist \( M>0 \) such that \( \abs{s_n}<M \) for all \( n\in\N_+ \). \par
  For any positive number \( \varepsilon \), there exist \( N_1,N_2\in\N_+ \) such that
  \[ \abs{s_mb_m-s_nb_n}<\frac{\varepsilon}{2}, \forall\,m>n>N_1;\quad \abs{b_m-b_n}<\frac{\varepsilon}{2M}, \forall\,m>n>N_2. \] 
  Let \( N=\max\{N_1,N_2\} \). For any \( m>n>N \), then for all \( m>n>N \),
  \begin{align*}
    \abs{S_m-S_n} & = \abs{s_mb_m-s_nb_n+\sum_{k=n}^{m-1}s_k(b_k-b_{k+1})} \\
    & \le \abs{s_mb_m-s_nb_n}+\sum_{k=n}^{m-1}\abs{s_k}\cdot\abs{b_k-b_{k+1}}
    < \frac{\varepsilon}{2}+M\sum_{k=n}^{m-1}\abs{b_k-b_{k+1}} \\
    & = \frac{\varepsilon}{2}+M\abs{b_{n+1}-b_m} \text{ (This is because $\{b_n\}$ is monotone)}\\
    &< \varepsilon.
  \end{align*}
  Therefore, \( \sum a_nb_n \) is a Cauchy series, and \( \sum a_nb_n \) converges.
\end{proof}



\begin{problem}[3-9]
Find the radius of convergence of each of the following power series: (a) $\displaystyle \sum n^3 z^n$, (b) $\displaystyle \sum \frac{2^n}{n!} z^n$, (c) $\displaystyle \sum \frac{2^n}{n^2} z^n$, (d) $\displaystyle \sum \frac{n^3}{3^n} z^n$.
\end{problem}

\begin{solution}
  (a) \[ R=\frac{1}{\limsup\sqrt[n]{n^3}} = \frac{1}{1} = 1. \]\par
  (b) \[ R=\frac{1}{\limsup\sqrt[n]{2^n/n!}} = \frac{1}{0} = \infty. \]\par
  (c) \[ R=\frac{1}{\limsup\sqrt[n]{2^n/n^2}} = \frac{1}{2\limsup n^{-2/n}} = \frac{1}{2}. \]\par
  (d) \[ R=\frac{1}{\limsup\sqrt[n]{n^3/3^n}} = \frac{1}{3\limsup n^{-3/n}} = \frac{1}{3}. \]\par
\end{solution}



\begin{problem}[3-10]
Assume that all coefficients of the power series
\[
\sum a_n z^n
\]
are integers, and infinitely many of them are nonzero. Prove that its radius of convergence is at most $1$.
\end{problem}

\begin{proof}
  It suffices to show that \( \limsup_{n\to\infty}\sqrt[n]{a_n}\ge 1 \). For any positive integer \( N \), there exists \( n>N \) such that \( a_n\ge 1 \) (since infinitely many of the coeffecients are nonzero), then the conclusion holds.
\end{proof}



\begin{problem}[3-14]
Let $\{s_n\}$ be a sequence of complex numbers, and define its arithmetic means by
\[
\sigma_n = \frac{s_0 + s_1 + \cdots + s_n}{n+1}, \qquad (n=0,1,2,\ldots).
\]
\begin{enumerate}
	\item[(a)] If $\lim_{n\to\infty} s_n = s$, prove that $\lim_{n\to\infty} \sigma_n = s$.
	\item[(b)] Construct a sequence $\{s_n\}$ such that $\lim_{n\to\infty} \sigma_n = 0$ but $\{s_n\}$ does not converge.
	\item[(c)] Can the following happen: $s_n>0$ for all $n$, $\lim_{n\to\infty} \sigma_n=0$, but $\limsup_{n\to\infty} s_n = \infty$?
	\item[(d)] For $n\ge 1$, let $a_n = s_n - s_{n-1}$. Prove that
	\[
	s_n - \sigma_n = \frac{1}{n+1}\sum_{k=1}^{n} k a_k.
	\]
	Assume further that $\lim_{n\to\infty}(n a_n)=0$ and that $\{\sigma_n\}$ converges. Prove that $\{s_n\}$ converges. (This is a partial converse of (a), with the additional assumption $n a_n\to 0$.)
	\item[(e)] Derive the conclusion in (d) under weaker hypotheses: assume $M<\infty$, $|n a_n|\le M$ for all $n$, and $\lim_{n\to\infty}\sigma_n=\sigma$. Prove $\lim_{n\to\infty}s_n=\sigma$ by following these steps:

	If $m<n$, then
	\[
	s_n-\sigma_n = \frac{m+1}{n-m}(\sigma_n-\sigma_m) + \frac{1}{n-m}\sum_{i=m+1}^{n}(s_n-s_i).
	\]
	For such $i$,
	\[
	|s_n-s_i| \le \frac{(n-i)M}{i+1} \le \frac{(n-m-1)M}{m+2}.
	\]
	Fix $\varepsilon>0$. For each $n$, choose a positive integer $m$ such that
	\[
	m \le \frac{n-\varepsilon}{1+\varepsilon} < m+1.
	\]
	Then
	\[
	\frac{m+1}{n-m} \le \frac{1}{\varepsilon},
	\]
	and $|s_n-s_i|<M\varepsilon$. Hence
	\[
	\limsup_{n\to\infty}|s_n-\sigma| \le M\varepsilon.
	\]
	Since $\varepsilon$ is arbitrary, conclude that $\lim_{n\to\infty}s_n=\sigma$.
\end{enumerate}
\end{problem}

\begin{proof}
  (a) For any positive number \( \varepsilon \), there exists \( N_1\in\N_+ \) such that \( \abs{s_n-s}<\varepsilon/3 \) for all \( n>N_1 \). Take \( N_2\in\N_+ \) such that \( s_0+s_1+\cdots+s_{N_1}<(\varepsilon/3)N_2 \). It's clear that there exists \( N_3\in\N_+ \) such that \( (N_1+1)s<N_3\varepsilon/3 \). \par
  Set \( N=\max \left\{N_1,N_2,N_3\right\} \). For any \( n>N \), 
  \begin{align*}
    \abs{\sigma_n-s}
    &\le\abs{\frac{s_0+s_1+\cdots+s_{N_1}}{n}}+\abs{\frac{s_{N_1+1}+s_{N_1+2}+\cdots+s_n}{n} - s} \\
    &\le\frac{\varepsilon}{3}+\abs{\frac{(n-N_1)s}{n}-s}+\frac{\abs{s_{N_1+1}-s}+\abs{s_{N_1+2}-s}+\cdots+\abs{s_n-s}}{n} \\
    &<\frac{\varepsilon}{3}+\frac{\varepsilon}{3}+\frac{\varepsilon}{3}
    =\varepsilon.
  \end{align*}\par
  Hence \( \lim_{n\to\infty}\sigma_n=s \). \par
  (b) Let \( s_n=(-1)^n \). Then \( \sigma_n=\frac{1}{n+1}\sum_{k=0}^{n}(-1)^k \). It's easy to verify that \( \lim_{n\to\infty}\sigma_n=0 \), but \( \{s_n\} \) does not converge. \par
  (c) Yes. Let
  \[ s_n=\begin{cases}
    k\,, & \text{if } n=2^k;\\
    2^{-n}, & \text{otherwise.}
  \end{cases} \]\par
  It's not hard to verify that \( s_n=\mathcal{O}(\log n) \), then \( \lim_{n\to\infty}\sigma_n=0 \). But by the definition of \( s_n \), we have \( \limsup_{n\to\infty}s_n=+\infty \).\par
  (d) Note that
  \begin{align*}
    \sum_{k=1}^{n} k a_k 
    & = \sum_{k=1}^{n} k (s_k - s_{k-1}) 
    = \sum_{k=1}^{n} k s_k - \sum_{k=1}^{n} k s_{k-1} 
    = n s_n + \sum_{k=1}^{n-1} k s_k - \sum_{k=0}^{n-1} (k+1) s_k \\
    & = n s_n - \sum_{k=0}^{n-1} s_k = (n+1)(s_n-\sigma_n).
  \end{align*}\par
  Suppose \( \lim_{n\to\infty}na_n=0 \) and \( \{\sigma_n\} \) converges. We only need to prove that \( \lim_{n\to\infty}(s_n-\sigma_n)=0 \). This is an immediate consequence of (a), with subsituting \( s_n \) to \( na_n \).\par
  (e) Assume \(M<\infty\), \(\abs{na_n}\le M\) for all \(n\), and \(\sigma_n\to\sigma\). 
  For \(m<n\), the identity in the statement gives
  \[ s_n-\sigma_n
    =\frac{m+1}{n-m}(\sigma_n-\sigma_m)
    +\frac{1}{n-m}\sum_{i=m+1}^{n}(s_n-s_i). \]
  For each \(i\in\{m+1,\dots,n\}\), \( s_n-s_i=\sum_{k=i+1}^{n}a_k \), hence
  \[ \abs{s_n-s_i}
    \le\sum_{k=i+1}^{n}\abs{a_k}
    \le\sum_{k=i+1}^{n}\frac{M}{k}
    \le\frac{(n-i)M}{i+1}
    \le\frac{(n-m-1)M}{m+2}. \]\par
  Fix \(\varepsilon>0\). For each \(n\), choose \(m\in\N_+\) such that
  \[ m\le\frac{n-\varepsilon}{1+\varepsilon}<m+1. \]
  Then \( n-m\ge\varepsilon(m+1)+\varepsilon>\varepsilon(m+1) \), so
  \[ \frac{m+1}{n-m}\le\frac1\varepsilon. \]
  Also, \( n-m-1<\varepsilon(m+2) \), therefore for \(i=m+1,\dots,n\),
  \[ \abs{s_n-s_i}
    \le\frac{(n-m-1)M}{m+2}
    <M\varepsilon. \]\par
    Taking absolute values in the decomposition of \(s_n-\sigma_n\), we get
  \[ \abs{s_n-\sigma_n}
    \le \frac{m+1}{n-m}\abs{\sigma_n-\sigma_m}
    +\frac{1}{n-m}\sum_{i=m+1}^{n}\abs{s_n-s_i}
    \le \frac1\varepsilon\abs{\sigma_n-\sigma_m}+M\varepsilon. \]
  Since \(m\to\infty\) as \(n\to\infty\) and \(\sigma_n\to\sigma\), we have \( \abs{\sigma_n-\sigma_m}\to0 \), hence \( \limsup_{n\to\infty}\abs{s_n-\sigma_n}\le M\varepsilon \). Finally, \( \abs{s_n-\sigma}\le\abs{s_n-\sigma_n}+\abs{\sigma_n-\sigma} \), and \(\sigma_n\to\sigma\) implies
  \[ \limsup_{n\to\infty}\abs{s_n-\sigma}\le M\varepsilon. \]
  Because \(\varepsilon>0\) is arbitrary, \(\lim_{n\to\infty}s_n=\sigma\).
\end{proof}



\end{document}